\documentclass[11pt]{article}
\usepackage[T1]{fontenc}
\usepackage[utf8]{inputenc}
\usepackage{amsmath,amssymb,amsthm,booktabs}
\usepackage[margin=1in]{geometry}
\usepackage{hyperref,xurl}
\newtheorem{theorem}{Theorem}
\title{Coupled cusp constraints for three specified\linebreak ideal right-angled four-facet types}
\author{Coxeter Polytopes Frontier Project}
\date{18 October 2026}
\begin{document}
\maketitle
\begin{abstract}
Consider a hypothetical entirely ideal right-angled hyperbolic
five-polytope whose facets are restricted to the regular 24-cell,
its facet double, and a particular sixteen-cell positive-carrier
polytope. Cusp incidences and Lorentz orthogonality force at least
30 sixteen-cell facets and 66 doubled-24-cell facets. At a distinguished
cusp of each sixteen-cell facet, a complete finite necessary-compatibility
test forces eight sixteen-cell incident facets and excludes one of two
possible box widths. One local labeling and an aggregate count assignment
survive. Neither is a geometric realization; the three-type spectrum
and the unrestricted ideal-five-dimensional problem remain open.
\end{abstract}

\section{The specified facet family}
Write $C$ for the regular ideal right-angled 24-cell and $D$ for its
reflected double across an octahedral facet. Let $T$ be the exact
sixteen-cell polytope in the separate positive-carrier construction:
starting with $D$, double across the successive facets with saved unit
normals
\[
 F=(1,0,1,-1,0),\quad G=(1,0,1,0,-1),\quad H=(1,0,1,1,0).
\]
The metric is $\operatorname{diag}(-1,1,1,1,1)$, and the construction's
halfspace-union argument proves that these are actual finite-volume
ideal right-angled polytopes. The full rational normal and vertex sets
are the immutable supplementary inputs specified in the reproduction
guide. No classification of all ideal right-angled four-polytopes is
being assumed here.

For each ideal vertex $v=(1,x)$, the six incident unit normals form
three opposite pairs. Each pair has sum $\lambda v$, with $\lambda>0$.
In upper-half-space coordinates sending $v$ to infinity, the horizontal
normal has unit Euclidean length. Thus the horospherical width is
proportional to $\lambda$, with one common scale for that cusp.
Complete exact reconstruction gives
\begin{center}
\begin{tabular}{ccccc}\toprule
Facet&$(1,1,1)$ cusps&$(1,1,2)$ cusps&$(1,2,2)$ cusps&$\kappa$\\\midrule
$C$&24&0&0&0\\
$D$&36&6&0&$-2$\\
$T$&195&49&14&6\\\bottomrule
\end{tabular}
\end{center}
Here $\kappa=3f_0-4f_3+24$. The three-facets of $T$ have vertex-count
histogram $6^{128},9^{48},14^{18},24^4$. Every cusp record retains
the widths, opposite-pair assignments and the sizes of its six actual
incident three-facets.

\section{Global incidence and orthogonality}
\begin{theorem}
Suppose an entirely ideal finite-volume right-angled hyperbolic
five-polytope has every facet isometric to $C,D$ or $T$. If their
numbers are $x,y,z$, respectively, then
\[
 y=3z-24,\qquad 6\mid z,\qquad z\ge30,\qquad y\ge66.
\]
\end{theorem}
\begin{proof}
The separately proved ideal-five-dimensional carrier identity
$\sum\kappa=48$ gives $-2y+6z=48$.
Every ambient cusp is a rectangular four-box. All its three-box facets
have one of the three shapes in the table. Every pair of side lengths
occurs in a triple, so all ratios are at most two; three distinct lengths
would give an unlisted triple. The possible four-box shapes are therefore
$1111,1112,1122,1222$, up to scale and permutation.
Write $a,b,c,d$ for their numbers. Their eight facets have respective
three-box type multiplicities
\[
 (8,0,0),\quad(2,6,0),\quad(0,4,4),\quad(2,0,6).
\]
Counting actual cusp occurrences yields
\begin{align*}
8a+2b+2d&=24x+36y+195z,\\
6b+4c&=6y+49z,\\
4c+6d&=14z.
\end{align*}
The difference of the last two gives $6(b-d)=6y+35z$, hence $6\mid z$.

Each of the 22 three-facets of $T$ with 14 or 24 vertices must meet
another $T$: neither $C$ nor $D$ has such a facet. These 22 neighbors
are distinct because the intersection of two convex facets is a single
face. The ridge-adjacency graph on the $z$ facets of type $T$ therefore
has at least $11z$ edges. It has no six-clique, since six mutually
orthogonal unit spacelike vectors cannot occur in Lorentz space of
positive index five.

An elementary weighted clique argument gives $e\le2z^2/5$ for a graph
without a six-clique. Maximize $\sum_{ij\in E}t_it_j$ for $t_i\ge0$,
$\sum t_i=1$, choosing a maximizer with minimal support. If two supported
vertices are nonadjacent, moving their combined weight to the one with
larger weighted neighbor sum does not reduce the objective and reduces
the support. Hence the support is a clique, of size at most five, and
the objective is at most $(1-1/5)/2=2/5$. Uniform weights give the claim.
Thus $11z\le2z^2/5$. Since $z>0$ by $y=3z-24$, this implies
$z\ge55/2$. Divisibility by six now gives $z\ge30$ and $y\ge66$.
\end{proof}

\section{A forced local eight-facet configuration}
The complete cusp table of $T$ has a unique vertex with three opposite
ridge-size pairs
\[
 (24,24),\quad(24,24),\quad(14,14),
\]
whose widths are $(4,4,8)$ in the saved normalization. All six neighboring
ambient facets at this vertex must be $T$. The ambient four-box has
widths $(t,4,4,8)$, where the original $T$ is perpendicular to the
$t$ direction. A facet perpendicular to an old width-four direction
has three-box widths $(t,4,8)$, so the allowed shape list forces
$t\in\{4,8\}$. The opposite of the original $T$ has shape $(1,1,2)$,
and must be $D$ or $T$.

\begin{theorem}
Under the same facet-spectrum restriction, the distinguished cusp has
$t=4$ and all eight incident ambient facets have type $T$.
\end{theorem}
\begin{proof}
For each of the four cases $(t,\text{opposite type})$, take every actual
cusp record of the required facet type. Allow all permutations and sign
flips of the three coordinate axes preserving widths up to a common
scale. On adjacent box facets, require equal vertex counts for their
common three-face. This is a necessary overdomain: it even allows cusp
symmetries which might not extend to the full four-polytope.
The complete finite check gives
\begin{center}
\begin{tabular}{ccc}\toprule
$t$&Opposite facet&Necessary matching\\\midrule
4&$D$&none\\
4&$T$&an explicit labeling\\
8&$D$&none\\
8&$T$&none\\\bottomrule
\end{tabular}
\end{center}
The exact domains and exhaustive negative checks are saved and replayed
by two separately implemented searches. The producer uses direct
backtracking; the standalone checker reconstructs the domains and uses
arc consistency followed by complete branching.
For $t=8$, there is also a direct table obstruction: every $(1,2,2)$
cusp of $T$ has nine-vertex three-facets on its four long-direction
sides. A required neighbor has this cusp shape, but the side facing
the original $T$ is in the long $t$ direction and has 24 vertices.
The $t=4,D$ exclusion uses the coupled adjacent-facet matching and is
not inferred from the width list alone.
\end{proof}

\section{Surviving counts and proof scope}
The nonnegative assignment
\[
 (x,y,z)=(361,66,30),\qquad (a,b,c,d)=(2051,241,105,0)
\]
satisfies the displayed incidence equations and the ideal-five-dimensional
Euler identities, with 457 facets and 2397 vertices. It also satisfies
the elementary lower bound $1+198+258=457$ imposed by a $T$ facet and
its distinct ridge-neighbor and cusp-opposite facets.
It is only an arithmetic witness. The surviving local labeling does
not establish full three-face isometries, a consistent boundary complex,
a Lorentz Gram realization or a chamber. Unlisted facet types are outside
both theorems.

Generation, separate-process replay and the standalone checker passed.
An independent full internal mathematical source/report/helper review
checked the universal arguments and finite-domain scope, without
re-executing them or reproving the inherited four-polytopes. No external
human review or end-to-end formal-kernel proof is claimed.
See \path{docs/ideal_ra5_cusp_compatibility_reproduce.md} for exact files,
commands, immutable input hashes and the internal/external lemma map.
The universal ideal-six-dimensional size and mixed-density consequences
are presented separately in
\emph{Necessary size and face-density constraints for ideal right-angled
six-polytopes}; they have no three-type facet-spectrum restriction.
\end{document}
